% Generated draft: requires mathematical and source review.
\documentclass[11pt]{article}
\usepackage[margin=25mm]{geometry}
\usepackage{fontspec}
\setmainfont{Latin Modern Roman}
\usepackage{amsmath,amssymb,mathtools,unicode-math}
\setmathfont{Latin Modern Math}
\usepackage{enumitem,needspace,xurl,hyperref}
\setlength{\parindent}{0pt}
\setlength{\parskip}{5pt}
\setlength{\emergencystretch}{3em}
\newcommand{\N}{\mathbb{N}}
\newcommand{\Z}{\mathbb{Z}}
\newcommand{\Q}{\mathbb{Q}}
\newcommand{\R}{\mathbb{R}}
\newcommand{\C}{\mathbb{C}}
\newcommand{\F}{\mathbb{F}}
\newcommand{\A}{\mathbb{A}}
\newcommand{\PP}{\mathbb{P}}
\newcommand{\E}{\mathbb{E}}
\newcommand{\eps}{\varepsilon}
\newcommand{\dd}{\,\mathrm{d}}
\begin{document}
\section*{UnsolvedMath ID: 30001705}
\textbf{Source problem code:} catalogue entry
\subsection{Definitions and mathematical statement}
A maximal torus \(T\) of a Lie group \(G\) is a compact connected abelian Lie subgroup maximal among such subgroups. Writing \(\mathfrak t\) for its Lie algebra, one may identify \(T\cong\mathfrak t/\Lambda\), where \(\Lambda=\ker(\exp\colon\mathfrak t\to T)\) is its period lattice. Once a Riemannian metric and a heat-equation convention are fixed, the heat kernel is the integral kernel of the corresponding heat evolution operator. A lattice-sum expression represents that kernel using a sum indexed by a lattice. The supplied record does not specify the ambient group, metric, lattice-sum formula, or heat-equation normalization; it also does not clarify whether the intended kernel is intrinsic to \(T\) or is an ambient-group heat kernel evaluated on \(T\).

For a smooth manifold \(M\), the smooth free loop space is \(LM=C^\infty(S^1,M)\). This defines one conventional meaning of loop space; the source does not specify the underlying manifold, whether loops are free or based, or the analytic completion to be used. Hodge theory concerns harmonic representatives of cohomology classes for a specified differential complex and analytic structure. Morse theory relates topology to critical points of a function or functional under suitable hypotheses. A Witten-type interpolation refers here to a deformation intended to connect these two descriptions. The differential complex, operators, analytic hypotheses, and limiting assertions for such an interpolation on loop space are not supplied.

\textbf{Source question.} Can the lattice-sum evaluation of the heat kernel on a maximal torus be related to Duistermaat--Heckman and Berline--Vergne localization? Can an analogous Witten-type interpolation between Hodge theory and Morse theory be constructed on loop space for a Lagrangian such as the energy?

These are construction and comparison questions. The source does not prescribe the mathematical meaning of the requested relation, the localization data, or the properties that a successful interpolation must satisfy. Energy is given as an example, without a specified formula, metric, domain, or boundary conditions. Thus the supplied evidence supports this qualified research question, but does not determine a unique fully quantified existence statement.
\subsection{Short English statement}
The problem asks whether a lattice-sum heat-kernel formula on a maximal torus admits an interpretation through Duistermaat--Heckman and Berline--Vergne localization, and whether a related Witten-type deformation can connect Hodge and Morse theories on loop space, for example using energy. The geometric and analytic setting and the required comparison properties remain unspecified.
\subsection{Sources}
\begin{itemize}
\item Supplied local UnsolvedMath registry record, permanent ID \(30001705\), source problem code \texttt{catalogue entry}, titled \emph{Loop-Space Localization and Witten-Type Interpolation}, in the Oberwolfach Reports Open Problems collection. The record assigns the proposed year 2011 and names no individual proposer.
\item \emph{Geometric Quantization in the Non-compact Setting} (2011), Oberwolfach Reports, pp.\ 425--521, DOI \texttt{10.4171/owr/2011/09}. Supplied source and literature-source link, labeled ``original Oberwolfach report'': \url{https://doi.org/10.4171/owr/2011/09}. The supplied background describes these reports as records of mathematical workshops and their open-problem discussions; no problem-specific page or speaker attribution is provided.
\item Supplied literature assessment, checked exactly \texttt{2026-08-21}: the record labels the problem open and describes the loop-space interpolation question as an obstruction to extending localization methods, with no resolving Hodge theory or general construction supplied. This is supplied dated curation, not an independently verified literature assessment. The background expressly requests independent verification of time-sensitive status and cited claims.
\item The registry labels the statement \texttt{corrected\_verified} and reports a manual semantic, typographic, and LaTeX review of its cleanup. These are supplied editorial claims, not independent verification of mathematical completeness or correctness.
\end{itemize}
\end{document}
